% ==================== SECTION LOI D'OHM ====================

% ========== VARIANTE A ==========
\element{loi-ohm}{
    \begin{question}{ohm-A-1}
        En supposant que toutes les variables sont correctement déclarées, en considérant le programme ci-dessous : compléter le tableau donnant les grandeurs électriques calculées par le programme.

        \textbf{Rappel :} $U = R \times I$ et $P = U \times I$

        \begin{minipage}{.55\textwidth}
            \lstinputlisting{sources/codes/loi_ohm_A.c}
        \end{minipage}\hfill
        \begin{minipage}{.4\textwidth}
            {
                \Large
                \begin{tabular}{|l|c|}
                    \hline
                    \textbf{Calcul} & \textbf{Résultat} \\\hline
                    Courant I1 (A) &  \\\hline
                    Puissance P1 (W) &  \\\hline
                    Courant I2 (A) &  \\\hline
                    Tension U2 (V) &  \\\hline
                \end{tabular}
            }
        \end{minipage}

        \evaluationProf
    \end{question}

    \begin{question}{ohm-A-2}
        Analyser le code de ce programme en apportant une critique de son organisation et de sa structure.

        \textit{Indices : répétitions, modularité, réutilisabilité...}

        \evaluationProfGrille{2}{8}
    \end{question}

    \begin{question}{ohm-A-3}
        Proposer une refactorisation du code en créant au minimum deux fonctions appropriées pour améliorer la structure du programme.

        \textbf{Écrire uniquement les prototypes et les définitions des fonctions (pas besoin de réécrire le main complet).}

        \evaluationProfGrille{2}{18}
    \end{question}
}

% ========== VARIANTE B ==========
\element{loi-ohm}{
    \begin{question}{ohm-B-1}
        En supposant que toutes les variables sont correctement déclarées, en considérant le programme ci-dessous : compléter le tableau donnant les grandeurs électriques calculées par le programme.

        \textbf{Rappel :} $U = R \times I$ et $P = U \times I$

        \begin{minipage}{.55\textwidth}
            \lstinputlisting{sources/codes/loi_ohm_B.c}
        \end{minipage}\hfill
        \begin{minipage}{.4\textwidth}
            {
                \Large
                \begin{tabular}{|l|c|}
                    \hline
                    \textbf{Calcul} & \textbf{Résultat} \\\hline
                    Courant I1 (A) &  \\\hline
                    Puissance P1 (W) &  \\\hline
                    Courant I2 (A) &  \\\hline
                    Tension U2 (V) &  \\\hline
                \end{tabular}
            }
        \end{minipage}

        \evaluationProf
    \end{question}

    \begin{question}{ohm-B-2}
        Analyser le code de ce programme en apportant une critique de son organisation et de sa structure.

        \textit{Indices : répétitions, modularité, réutilisabilité...}

        \evaluationProfGrille{2}{8}
    \end{question}

    \begin{question}{ohm-B-3}
        Proposer une refactorisation du code en créant au minimum deux fonctions appropriées pour améliorer la structure du programme.

        \textbf{Écrire uniquement les prototypes et les définitions des fonctions (pas besoin de réécrire le main complet).}

        \evaluationProfGrille{2}{18}
    \end{question}
}

% ========== VARIANTE C ==========
\element{loi-ohm}{
    \begin{question}{ohm-C-1}
        En supposant que toutes les variables sont correctement déclarées, en considérant le programme ci-dessous : compléter le tableau donnant les grandeurs électriques calculées par le programme.

        \textbf{Rappel :} $U = R \times I$ et $P = U \times I$

        \begin{minipage}{.55\textwidth}
            \lstinputlisting{sources/codes/loi_ohm_C.c}
        \end{minipage}\hfill
        \begin{minipage}{.4\textwidth}
            {
                \Large
                \begin{tabular}{|l|c|}
                    \hline
                    \textbf{Calcul} & \textbf{Résultat} \\\hline
                    Courant I1 (A) &  \\\hline
                    Puissance P1 (W) &  \\\hline
                    Courant I2 (A) &  \\\hline
                    Tension U2 (V) &  \\\hline
                \end{tabular}
            }
        \end{minipage}

        \evaluationProf
    \end{question}

    \begin{question}{ohm-C-2}
        Analyser le code de ce programme en apportant une critique de son organisation et de sa structure.

        \textit{Indices : répétitions, modularité, réutilisabilité...}

        \evaluationProfGrille{2}{8}
    \end{question}

    \begin{question}{ohm-C-3}
        Proposer deux fonctions : une pour calculer la tension $U$ et une pour calculer le courant $I$, afin d'améliorer la structure du programme.

        \textbf{Écrire uniquement les prototypes et les définitions des fonctions (pas besoin de réécrire le main complet).}

        \evaluationProfGrille{2}{18}
    \end{question}
}
